% =============================================================================
% Project : THE PHILOSOPHY OF VIABILITY
% Book    : Book I — Ontology of Viability
% Book ID : VS.B01
% Part    : Part I — Identity and Manifestation
% Part ID : VS.B01.P01
% Chapter : Chapter 08 — Comparative Philosophical Examination
% Chapter ID : VS.B01.P01.C08
%
% File :
%   2_MAINMATTER/PART_I_IDENTITY_AND_MANIFESTATION/
%   CHAPTER_08_COMPARATIVE_PHILOSOPHICAL_EXAMINATION/
%   sections/section_15_averroes.tex
%
% Section :
%   Section 15 — Averroes (Ibn Rushd)
%
% Structural Identifier :
%   VS.B01.P01.C08.S15
%
% Purpose :
%   To examine Averroes's account of being, substance, matter,
%   form, essence, unity, potentiality, actuality, causality,
%   and change, and to test whether this account can replace
%   or exhaust the ontological role assigned to Identity in
%   Book I.
%
% Method :
%   Historical reconstruction, textual differentiation,
%   bilateral criticism, explanatory-cost assessment,
%   and a substitution test.
%
% Dependencies :
%   Section 01 — Method and Criteria of Comparison.
%   Section 13 — Avicenna.
%   Section 14 — Suhrawardi.
%
% Assets :
%   TL01 — Averroes's Principal Claims and Comparative Limits.
%   DG01 — Averroes's Account of Substance and the Book I
%          Relation between Being and Identity.
%
% Bibliographical Basis :
%   AristotleMetaphysics
%   AristotlePhysics
%   AverroesEpitomeMetaphysicsArnzen2010
%   AverroesLongCommentaryMetaphysicsBouyges1990
%   AverroesMetaphysicsBookLambdaGenequand1984
%   AverroesIncoherenceVanDenBergh1954
%   AverroesDecisiveTreatise
%   BenAhmedPasnauIbnRushdSEP2025
%   DiGiovanniIndividuationMatterAverroes2007
%   DiGiovanniSubstantialFormAverroes2011
%   BeloEssenceExistenceAvicennaAverroes2009
%   BertolacciAvicennaAverroesMetaphysics2007
%   ShayeganAccidentalityExistence2017
%   TaylorAverroesUnityOfTruth2000
%
% =============================================================================


\section{Averroes (Ibn Rushd)}
\label{sec:ch08_s15}


% =============================================================================
% SS01 — SCOPE AND COMPARATIVE QUESTION
% =============================================================================

\subsection{Scope and Comparative Question}
\label{subsec:ch08_s15_ss01}

% Structural ID: VS.B01.P01.C08.S15.SS01.P001
\phantomsection
\label{para:VS.B01.P01.C08.S15.SS01.P001}

The present examination concerns Ibn Rushd's metaphysical account of
being and the constitution of sensible substances. Its central
question is not whether Averroes anticipated the terminology of
Book I, but whether the explanatory principles of his Aristotelian
ontology leave a distinct ontological problem that can be formulated
as the problem of Identity. The comparison therefore begins with
Averroes's own concepts and argumentative purposes rather than
translating them prematurely into the vocabulary of the Philosophy
of Viability.

% Structural ID: VS.B01.P01.C08.S15.SS01.P002
\phantomsection
\label{para:VS.B01.P01.C08.S15.SS01.P002}

Averroes offers a powerful account of determinate sensible being
through substance, matter, form, essence, potentiality, actuality,
and causal order. These concepts explain how a sensible individual
is constituted, how it possesses a determinate nature, and how it
undergoes accidental and substantial change. They provide a serious
challenge to any ontology that introduces Identity without showing
what explanatory work it performs beyond the constitution and
determination of a substance.

% Structural ID: VS.B01.P01.C08.S15.SS01.P003
\phantomsection
\label{para:VS.B01.P01.C08.S15.SS01.P003}

The central comparative thesis is that Averroes's account does not,
without further argument, settle the question of Identity as
formulated in Book I. Substance, matter, form, essence, individuation,
unity, and persistence perform distinguishable explanatory functions.
The fact that these principles jointly explain a determinate existent
does not establish that any one of them, or their unexamined sum,
is equivalent to Identity. Conversely, Book I must demonstrate that
its concept of Identity is not an empty name for what Averroes
already explains through substantial determination.

% Structural ID: VS.B01.P01.C08.S15.SS01.P004
\phantomsection
\label{para:VS.B01.P01.C08.S15.SS01.P004}

The comparison is restricted to the ontological problems relevant
to Part I. Questions concerning the structure of human cognition,
the unity of the material and agent intellects, the conditions of
demonstrative knowledge, and the relation between philosophy and
religion are not developed as foundations of this section. They
belong to other inquiries and cannot be used to make the existence
or identity of an existent dependent upon a knower. The metaphysical
status of the agent intellect is likewise not treated as a
definition of Light in Book I.


% =============================================================================
% SS02 — HISTORICAL AND TEXTUAL LOCATION
% =============================================================================

\subsection{Historical and Textual Location}
\label{subsec:ch08_s15_ss02}

% Structural ID: VS.B01.P01.C08.S15.SS02.P001
\phantomsection
\label{para:VS.B01.P01.C08.S15.SS02.P001}

Abū al-Walīd Muḥammad ibn Aḥmad Ibn Rushd (1126–1198), commonly
known in Latinized form as Averroes, produced a large body of
commentaries and independent treatises. His relation to Aristotle
cannot be adequately described as passive repetition. His writings
include different literary forms, among them epitomes, paraphrases,
literal commentaries, and treatises directed toward particular
problems. These forms differ in their relation to the Aristotelian
text and in the degree to which Averroes reorganizes, interprets,
or criticizes its arguments
\citep{BenAhmedPasnauIbnRushdSEP2025}.

% Structural ID: VS.B01.P01.C08.S15.SS02.P002
\phantomsection
\label{para:VS.B01.P01.C08.S15.SS02.P002}

The *Epitome of the Metaphysics* must therefore be distinguished
from the Middle Commentary and the Long Commentary. It does not
simply reproduce the order of Aristotle's treatise. It selects and
reorganizes doctrines according to Averroes's conception of the
discipline. The Long Commentary, by contrast, follows Aristotle's
text more closely and belongs to a later stage of Averroes's
philosophical work. The distinction matters because a statement
found in an early epitome cannot automatically be treated as the
final formulation of a position developed in a later commentary
\citep{AverroesEpitomeMetaphysicsArnzen2010,
AverroesLongCommentaryMetaphysicsBouyges1990}.

% Structural ID: VS.B01.P01.C08.S15.SS02.P003
\phantomsection
\label{para:VS.B01.P01.C08.S15.SS02.P003}

The textual history of the *Epitome* itself requires caution.
Its extant transmission reflects more than one stage of revision,
and the work as transmitted is incomplete. Its structure and
selection of material are consequently relevant to interpretation.
The present section does not treat the *Epitome*, the Middle
Commentary, the Long Commentary, and the *Tahāfut al-Tahāfut*
as interchangeable witnesses to one unchanging formulation.
Each is used according to its genre, date, and subject
\citep{AverroesEpitomeMetaphysicsArnzen2010,
BenAhmedPasnauIbnRushdSEP2025}.

% Structural ID: VS.B01.P01.C08.S15.SS02.P004
\phantomsection
\label{para:VS.B01.P01.C08.S15.SS02.P004}

The inherited comparative manuscript in
*The Philosophy of Viability in Dialogue*, Chapter 41, contains
useful historical material but cannot serve as the final
interpretation of Averroes. In particular, its presentation of
a doctrine of two mutually opposed truths is not retained.
The familiar description of Averroes as a proponent of
contradictory philosophical and religious truths is misleading:
his discussion of the relation between truth and philosophical
inquiry must be read in the context of his account of the unity
of truth and the different forms of argument
\citep{AverroesDecisiveTreatise,
TaylorAverroesUnityOfTruth2000}.

% Structural ID: VS.B01.P01.C08.S15.SS02.P005
\phantomsection
\label{para:VS.B01.P01.C08.S15.SS02.P005}

The primary textual basis of the present examination is therefore
the Aristotelian metaphysical corpus as interpreted by Averroes,
especially the *Epitome* and the Long Commentary. The
*Tahāfut al-Tahāfut* is used where the discussion of causality,
generation, and the eternity of the world requires it. The
*Decisive Treatise* is not used as a substitute for the
metaphysical commentaries. It is relevant only to the limited
historical correction made above.


% =============================================================================
% SS03 — BEING AND THE SUBJECT OF METAPHYSICS
% =============================================================================

\subsection{Being and the Subject of Metaphysics}
\label{subsec:ch08_s15_ss03}

% Structural ID: VS.B01.P01.C08.S15.SS03.P001
\phantomsection
\label{para:VS.B01.P01.C08.S15.SS03.P001}

Averroes approaches metaphysics as a universal theoretical science
concerned with being insofar as it is being and with the principles
that belong to being universally. Its subject is not one
departmental kind of existent, such as a natural body or a
mathematical magnitude. It investigates what cannot be adequately
examined by a particular science because those matters concern
existents across the divisions of the special sciences. This
conception follows the Aristotelian determination of first
philosophy while giving it a systematic role in the organization
of theoretical inquiry
\citep{AristotleMetaphysics,
AverroesEpitomeMetaphysicsArnzen2010}.

% Structural ID: VS.B01.P01.C08.S15.SS03.P002
\phantomsection
\label{para:VS.B01.P01.C08.S15.SS03.P002}

Being is not treated as a genus whose instances differ only by
additional specific features. The different senses and modes
of being must be considered through their ordered relations,
with substance occupying a primary position. The primacy of
substance is not a claim that every existent is a substance
in exactly the same sense. It indicates that other categories,
including accidents, depend upon substance in a way that
substance does not depend upon them. Metaphysics must therefore
examine the structure of substance if it is to understand the
universal order of being
\citep{AristotleMetaphysics,
BenAhmedPasnauIbnRushdSEP2025}.

% Structural ID: VS.B01.P01.C08.S15.SS03.P003
\phantomsection
\label{para:VS.B01.P01.C08.S15.SS03.P003}

Averroes distinguishes what is most familiar to us from what
is most fundamental in the constitution of sensible substances.
Concrete individuals, such as animals and stones, are the
familiar substances encountered in ordinary experience.
Yet these individuals are composites whose principles require
metaphysical examination. Matter and form are therefore not
merely two names for observable features of an object; they
are principles through which the composite is constituted
\citep{AverroesEpitomeMetaphysicsArnzen2010,
AverroesLongCommentaryMetaphysicsBouyges1990,
BenAhmedPasnauIbnRushdSEP2025}.

% Structural ID: VS.B01.P01.C08.S15.SS03.P004
\phantomsection
\label{para:VS.B01.P01.C08.S15.SS03.P004}

The comparative consequence is important. Averroes's primary
metaphysical question concerns the existent and the principles
that account for its substantial constitution. Book I asks
a related but non-identical question concerning the ontological
relation between Being and Identity. The Aristotelian primacy
of substance does not itself establish that substance is
identical with Being as treated in Book I, nor does it establish
that the form or essence of a substance is identical with
Identity. Those identifications would require additional
arguments.


% =============================================================================
% SS04 — SUBSTANCE AND THE COMPOSITE INDIVIDUAL
% =============================================================================

\subsection{Substance and the Composite Individual}
\label{subsec:ch08_s15_ss04}

% Structural ID: VS.B01.P01.C08.S15.SS04.P001
\phantomsection
\label{para:VS.B01.P01.C08.S15.SS04.P001}

Substance is the primary mode of being in Averroes's account
of sensible reality. The concrete individual is not merely
a collection of qualities, nor is it a neutral subject to
which all determinations are externally attached. It is a
substantial existent whose determinate nature depends upon
the principles through which it is constituted. The account
of substance consequently gives a central place to the
relation between matter and form
\citep{AverroesEpitomeMetaphysicsArnzen2010,
BenAhmedPasnauIbnRushdSEP2025}.

% Structural ID: VS.B01.P01.C08.S15.SS04.P002
\phantomsection
\label{para:VS.B01.P01.C08.S15.SS04.P002}

Accidents are dependent features of a substance. A substance
may acquire or lose accidental determinations while remaining
the subject of those changes. The dependence of an accident
upon a substance is not equivalent to the dependence of
matter upon form. These relations occupy different explanatory
roles: an accident belongs to a substance as a dependent
determination, whereas matter and form are principles within
the account of substantial constitution
\citep{AverroesEpitomeMetaphysicsArnzen2010,
AverroesLongCommentaryMetaphysicsBouyges1990}.

% Structural ID: VS.B01.P01.C08.S15.SS04.P003
\phantomsection
\label{para:VS.B01.P01.C08.S15.SS04.P003}

The composite substance is not adequately explained by
treating matter as a complete individual that merely receives
external properties. Matter is potential with respect to
forms, while form accounts for the actuality and determinate
nature of the composite. The substance is thus intelligible
through the relation of its principles rather than through
the isolated consideration of a bare substratum
\citep{AverroesLongCommentaryMetaphysicsBouyges1990,
BenAhmedPasnauIbnRushdSEP2025}.

% Structural ID: VS.B01.P01.C08.S15.SS04.P004
\phantomsection
\label{para:VS.B01.P01.C08.S15.SS04.P004}

For Book I, this analysis establishes an important comparative
limit. Identity cannot be introduced as another accidental
feature added to a substance, because the foundational
commitments of Book I do not treat Identity as an additional
property. Nor can Identity simply be identified with the
composite substance. Substance explains a mode of ontological
constitution; the distinct question of Identity concerns
the ontological status of that which is identity-bearing.
The relation between these questions requires argument,
not terminological substitution.


% =============================================================================
% SS05 — PRIME MATTER AND POTENTIALITY
% =============================================================================

\subsection{Prime Matter and Potentiality}
\label{subsec:ch08_s15_ss05}

% Structural ID: VS.B01.P01.C08.S15.SS05.P001
\phantomsection
\label{para:VS.B01.P01.C08.S15.SS05.P001}

Prime matter is not a fully actual sensible object existing
independently of every form. It is considered through its
potentiality to receive forms and to serve as the underlying
principle in generation and change. Its explanatory function
is to account for the material capacity through which
different substantial forms can be actualized. Matter is
therefore not equivalent to a determinate body, even though
it is indispensable to the account of corporeal substances
\citep{AverroesEpitomeMetaphysicsArnzen2010,
AverroesLongCommentaryMetaphysicsBouyges1990}.

% Structural ID: VS.B01.P01.C08.S15.SS05.P002
\phantomsection
\label{para:VS.B01.P01.C08.S15.SS05.P002}

The status of prime matter is philosophically difficult.
Averroes can treat matter as a substance and as the enduring
principle underlying generation, while also recognizing
that matter considered without form does not possess the
actuality of a concrete sensible existent. The attempt to
describe matter as pure potentiality therefore generates
a tension between its role as a real principle of change
and the difficulty of assigning it an independent actuality.
This tension must not be removed by turning prime matter
into a fully determinate entity
\citep{AverroesLongCommentaryMetaphysicsBouyges1990,
BenAhmedPasnauIbnRushdSEP2025}.

% Structural ID: VS.B01.P01.C08.S15.SS05.P003
\phantomsection
\label{para:VS.B01.P01.C08.S15.SS05.P003}

In some formulations, prime matter is described as homogeneous
and numerically one across transient corporeal bodies.
This formulation concerns the common material principle
underlying generation; it does not establish that the
generated bodies share one numerical identity. The unity
attributed to prime matter must be distinguished from the
numerical unity of a concrete substance and from the
ontological question of whether one or more identities
belong to an existent
\citep{AverroesLongCommentaryMetaphysicsBouyges1990,
BenAhmedPasnauIbnRushdSEP2025}.

% Structural ID: VS.B01.P01.C08.S15.SS05.P004
\phantomsection
\label{para:VS.B01.P01.C08.S15.SS05.P004}

The comparison with Book I prevents a materialist reduction
of Identity. Material constitution can explain the
corporeal conditions under which a substance is generated,
maintained, or transformed. It cannot be renamed Identity
without collapsing the distinction between the material
principle of a substance and the ontological status of
the identity-bearing existent. Book I therefore treats
material constitution as a distinct explanatory dimension.


% =============================================================================
% SS06 — FORM AND SUBSTANTIAL DETERMINATION
% =============================================================================

\subsection{Form and Substantial Determination}
\label{subsec:ch08_s15_ss06}

% Structural ID: VS.B01.P01.C08.S15.SS06.P001
\phantomsection
\label{para:VS.B01.P01.C08.S15.SS06.P001}

Form is the principle through which matter is actualized
as a determinate sensible substance. It accounts for the
specific nature of the composite and explains why a
substance is the kind of thing that it is. In this respect,
form is not simply an observable shape or an external
configuration. Substantial form belongs to the explanatory
constitution of the substance itself
\citep{AristotleMetaphysics,
AverroesLongCommentaryMetaphysicsBouyges1990,
DiGiovanniSubstantialFormAverroes2011}.

% Structural ID: VS.B01.P01.C08.S15.SS06.P002
\phantomsection
\label{para:VS.B01.P01.C08.S15.SS06.P002}

Averroes gives form a primary role in the explanation of
substance. The form is not merely one component placed
beside matter without explanatory asymmetry. It is the
principle through which the composite becomes a determinate
substance. His account of essence or quiddity consequently
assigns a decisive role to form rather than treating the
essence of a sensible thing as the simple sum of its matter
and form. This interpretation is supported by his discussion
of substantial form in the Long Commentary
\citep{AverroesLongCommentaryMetaphysicsBouyges1990,
DiGiovanniSubstantialFormAverroes2011}.

% Structural ID: VS.B01.P01.C08.S15.SS06.P003
\phantomsection
\label{para:VS.B01.P01.C08.S15.SS06.P003}

The theory of substantial form must also be read historically.
Averroes's earlier discussions of generation sometimes
assign a role to a separate giver of forms, whereas later
formulations give greater explanatory weight to the
potentialities of prime matter and the activity of natural
agents. The development is significant because it shows
that his account of generation was not fixed in every
detail throughout his career. A comparison that ignores
this development would attribute to the mature position
claims that belong to an earlier explanatory framework
\citep{BenAhmedPasnauIbnRushdSEP2025}.

% Structural ID: VS.B01.P01.C08.S15.SS06.P004
\phantomsection
\label{para:VS.B01.P01.C08.S15.SS06.P004}

The priority of form in explaining substantial determination
does not establish that form is identical with Identity.
Form explains the determinate nature and essence of a
substance. Identity, as treated in Book I, is not defined
as the formal cause of a thing, nor is it introduced as
a substitute for the substantial form. The comparative
question is whether the ontological role assigned to
Identity can be stated without reproducing the explanatory
function already performed by form.


% =============================================================================
% SS07 — ESSENCE AND QUIDDITY
% =============================================================================

\subsection{Essence, Quiddity, and the Determinate Thing}
\label{subsec:ch08_s15_ss07}

% Structural ID: VS.B01.P01.C08.S15.SS07.P001
\phantomsection
\label{para:VS.B01.P01.C08.S15.SS07.P001}

Essence and quiddity concern what a thing is. They belong
to the analysis of the determinate nature of an existent
and must be distinguished from the bare affirmation that
the existent is. The terms *dhāt* and *māhiyya* occur in
different argumentative contexts and should not be treated
as interchangeable in every passage. Their translation
and interpretation require attention to the specific
text and the role played by the term in the argument
\citep{AverroesEpitomeMetaphysicsArnzen2010,
AverroesLongCommentaryMetaphysicsBouyges1990}.

% Structural ID: VS.B01.P01.C08.S15.SS07.P002
\phantomsection
\label{para:VS.B01.P01.C08.S15.SS07.P002}

Averroes's account gives substantial form a decisive role
in explaining the essence of a sensible substance.
The essence is not simply the material substratum, nor
is it an inventory of all the accidents that happen to
belong to an individual. It is explained through the
form by which the composite is the determinate substance
that it is. This position creates a direct challenge
for any account that assigns Identity a role close to
whatness without specifying the difference between
Identity and essence
\citep{AverroesLongCommentaryMetaphysicsBouyges1990,
DiGiovanniSubstantialFormAverroes2011}.

% Structural ID: VS.B01.P01.C08.S15.SS07.P003
\phantomsection
\label{para:VS.B01.P01.C08.S15.SS07.P003}

Book I cannot resolve this challenge by renaming essence
as Identity. Essence and Identity may concern overlapping
questions, but conceptual proximity does not establish
conceptual equivalence. Essence explains what a substance
is through its determinate nature; Identity is examined
as a distinct ontological concern that is inseparable
from Being without being reduced to any single determination.
If the two are to be related, the relation must be
argued rather than assumed.

% Structural ID: VS.B01.P01.C08.S15.SS07.P004
\phantomsection
\label{para:VS.B01.P01.C08.S15.SS07.P004}

The comparison also exposes a positive requirement for
Book I. It is insufficient to state only that Identity
is not essence, form, or matter. The theory must explain
what ontological question is answered by Identity and
why that question is not already answered by the account
of substantial constitution. This requirement becomes
the central test in the subsequent examination.


% =============================================================================
% SS08 — EXISTENCE AND THE CRITIQUE OF AVICENNA
% =============================================================================

\subsection{Existence and the Critique of Avicenna}
\label{subsec:ch08_s15_ss08}

% Structural ID: VS.B01.P01.C08.S15.SS08.P001
\phantomsection
\label{para:VS.B01.P01.C08.S15.SS08.P001}

The relation between essence and existence provides the
most direct point of contact between Averroes and the
preceding examination of Avicenna in
\StructuralRef{sec:ch08_s13}. Avicenna distinguishes the
quiddity of a contingent existent from the fact that it
exists. Considered in itself, the quiddity does not
entail existence; the existence of a contingent existent
therefore requires a causal explanation. This distinction
has modal and explanatory consequences and cannot be
reduced to a merely verbal difference
\citep{AvicennaMetaphysics,
BenAhmedPasnauIbnRushdSEP2025}.

% Structural ID: VS.B01.P01.C08.S15.SS08.P002
\phantomsection
\label{para:VS.B01.P01.C08.S15.SS08.P002}

Averroes criticizes important parts of Avicenna's treatment
of existence and its relation to essence. His criticism
is directed against the way existence is treated as a
common attribute or accident of the existent. The
Aristotelian alternative gives priority to the analysis
of being through its modes and categories, especially
substance, rather than treating existence as a feature
added uniformly to an already intelligible essence.
The disagreement concerns the ontological and explanatory
status of existence, not simply the use of two different
words
\citep{AverroesLongCommentaryMetaphysicsBouyges1990,
BertolacciAvicennaAverroesMetaphysics2007}.

% Structural ID: VS.B01.P01.C08.S15.SS08.P003
\phantomsection
\label{para:VS.B01.P01.C08.S15.SS08.P003}

The Avicennian position must nevertheless be represented
with care. The distinction between essence and existence
does not mean that a contingent essence first exists as
a separate entity and only afterwards receives existence
as an additional object. Avicenna's analysis concerns
the essence considered in itself and the actuality of
its existence, with a corresponding distinction in
causal and modal explanation. The precise ontological
status of the distinction remains a matter of scholarly
interpretation. Averroes's criticism therefore cannot
be treated as a decisive refutation merely by repeating
the description of existence as an accident
\citep{BeloEssenceExistenceAvicennaAverroes2009,
ShayeganAccidentalityExistence2017}.

% Structural ID: VS.B01.P01.C08.S15.SS08.P004
\phantomsection
\label{para:VS.B01.P01.C08.SS08.P004}

The modern discussion includes a significant counter-reading:
Shayegan argues that Averroes's criticism rests on a
misunderstanding of Avicenna's distinction and of the
relevant Aristotelian terminology. This interpretation
must itself be assessed critically; it is not adopted
as an unquestionable verdict. Its importance is that
it prevents the comparison from presenting Averroes's
criticism as though the Avicennian position were a
simple theory of two independently existing components
called essence and existence
\citep{ShayeganAccidentalityExistence2017,
BeloEssenceExistenceAvicennaAverroes2009}.

% Structural ID: VS.B01.P01.C08.S15.SS08.P005
\phantomsection
\label{para:VS.B01.P01.C08.S15.SS08.P005}

The comparison with Book I is not a choice between
Averroes and Avicenna on the meaning of existence.
The foundational relation between Being and Identity
in Book I is analytically distinct from both the
Avicennian distinction between quiddity and existence
and the Averroist critique of existence as a common
attribute. Neither the word *existence* nor the word
*essence* can be substituted for Identity without
changing the conceptual question under examination.


% =============================================================================
% SS09 — UNITY, THE ONE, AND THE MANY
% =============================================================================

\subsection{Unity, the One, and the Many}
\label{subsec:ch08_s15_ss09}

% Structural ID: VS.B01.P01.C08.S15.SS09.P001
\phantomsection
\label{para:VS.B01.P01.C08.S15.SS09.P001}

The treatment of unity and multiplicity forms a significant
part of Averroes's engagement with Aristotle's *Metaphysics*.
The *Epitome* reorganizes the material concerning the one
and the many and treats their relations as a central
metaphysical problem. The question is not limited to
counting separate objects. It concerns the ways in
which unity and multiplicity are predicated and how
they relate to the different modes of being
\citep{AristotleMetaphysics,
AverroesEpitomeMetaphysicsArnzen2010}.

% Structural ID: VS.B01.P01.C08.S15.SS09.P002
\phantomsection
\label{para:VS.B01.P01.C08.S15.SS09.P002}

Averroes's discussion belongs to an Aristotelian framework
in which unity is not simply an additional material
component of a substance. Unity and multiplicity must
be understood through their relation to being and to
the determinate structures of things. The different
senses of unity cannot be collapsed into one undifferentiated
notion, and the analysis of unity must distinguish
the unity of a substance from the unity of an aggregate
or a collection.

% Structural ID: VS.B01.P01.C08.S15.SS09.P003
\phantomsection
\label{para:VS.B01.P01.C08.S15.SS09.P003}

This analysis is relevant to Book I because the
ontological unity of a being cannot be equated with
Identity. Unity, singularity, numerical identity,
aggregate unity, and compositional unity answer
different questions. A single substance may possess
a unity that is not adequately described by treating
it as a mere aggregate; this does not establish that
unity itself is the ontological status of Identity.
The distinctions developed in Chapter 03 remain
necessary for the comparison.

% Structural ID: VS.B01.P01.C08.S15.SS09.P004
\phantomsection
\label{para:VS.B01.P01.C08.S15.SS09.P004}

The comparison leaves open whether the Aristotelian
analysis of the one and the many can account for every
form of plurality relevant to Book I. In particular,
the relation between the unity of a concrete substance
and the possible plurality of identities is not settled
by the classification of unity alone. Book I therefore
retains the question as an open ontological problem
rather than deriving a universal answer from the
Aristotelian framework.


% =============================================================================
% SS10 — INDIVIDUATION AND MULTIPLICATION
% =============================================================================

\subsection{Individuation, Multiplication, and Numerical Plurality}
\label{subsec:ch08_s15_ss10}

% Structural ID: VS.B01.P01.C08.S15.SS10.P001
\phantomsection
\label{para:VS.B01.P01.C08.S15.SS10.P001}

The question of individuation must be separated from
the question of substantial form. Averroes treats forms
as individual in their existence, while universals
do not exist as universals outside the soul. This
position means that form does not require an additional
principle that first turns an otherwise universal
form into an individual form. The account of material
multiplicity must therefore be distinguished from
the account of the individuality of form
\citep{AverroesEpitomeMetaphysicsArnzen2010,
BenAhmedPasnauIbnRushdSEP2025}.

% Structural ID: VS.B01.P01.C08.S15.SS10.P002
\phantomsection
\label{para:VS.B01.P01.C08.S15.SS10.P002}

Extended matter plays a role in the multiplication
of material forms. In the relevant interpretation,
matter is better described as a principle of multiplication
than as a universal principle of individuation. It
provides the material conditions under which more
than one form can be instantiated in the corporeal
realm. The distinction is important: multiplication
explains how material plurality is possible, whereas
individuation concerns what makes an existent this
individual. These questions are related but not
interchangeable
\citep{DiGiovanniIndividuationMatterAverroes2007,
BenAhmedPasnauIbnRushdSEP2025}.

% Structural ID: VS.B01.P01.C08.S15.SS10.P003
\phantomsection
\label{para:VS.B01.P01.C08.S15.SS10.P003}

The account does not establish that every substance
has exactly one Identity in the sense adopted by
Book I. Nor does the fact that forms are individual
settle whether plurality of identities is possible
at different ontological levels or under different
respects. The Aristotelian analysis of matter and
form explains important conditions of numerical
multiplicity, but it does not supply a complete
criterion of Identity as formulated in this book.

% Structural ID: VS.B01.P01.C08.S15.SS10.P004
\phantomsection
\label{para:VS.B01.P01.C08.S15.SS10.P004}

The distinction is especially important for the
current foundations. Identity is not equated with
Individuation, Numerical Identity, Singularity, or
the material conditions of multiplication. The
discussion of Averroes thus provides a concrete
case in which the explanatory role of individuation
can be recognized without treating it as an exhaustive
account of Identity.


% =============================================================================
% SS11 — POTENTIALITY, ACTUALITY, AND CHANGE
% =============================================================================

\subsection{Potentiality, Actuality, and Natural Change}
\label{subsec:ch08_s15_ss11}

% Structural ID: VS.B01.P01.C08.S15.SS11.P001
\phantomsection
\label{para:VS.B01.P01.C08.S15.SS11.P001}

Potentiality and actuality are central to Averroes's
analysis of change. Potentiality concerns the capacity
of a subject to receive or undergo a determination;
actuality concerns the realization of that capacity.
The relation is not a simple chronological sequence
in which every potentiality must first exist as an
independent object and then become actual. It is
an explanatory distinction within the analysis of
natural change and the principles of sensible
substances
\citep{AristotleMetaphysics,
AristotlePhysics,
AverroesEpitomeMetaphysicsArnzen2010}.

% Structural ID: VS.B01.P01.C08.S15.SS11.P002
\phantomsection
\label{para:VS.B01.P01.C08.S15.SS11.P002}

Prime matter is potential with respect to substantial
forms, while form accounts for the actuality of
the determinate composite. The analysis therefore
connects the material principle of change with the
formal principle through which a substance is
actualized. This account also requires distinctions
between different kinds of potentiality, since
the capacities of material subjects, natural agents,
and celestial bodies are not all described in
the same way
\citep{AverroesEpitomeMetaphysicsArnzen2010,
AverroesLongCommentaryMetaphysicsBouyges1990}.

% Structural ID: VS.B01.P01.C08.S15.SS11.P003
\phantomsection
\label{para:VS.B01.P01.C08.S15.SS11.P003}

The distinction between potentiality and actuality
must not be imported into Book I as a definition
of Viability. The possibility of receiving a form,
the capacity for change, and the actualization of
a natural power do not by themselves explain
Viability. Viability cannot be reduced to potentiality,
actuality, survival, persistence, successful
functioning, or adaptation. The Aristotelian
distinction is useful for examining change, but
it is not a substitute for the ontological concept
that Book I is developing.

% Structural ID: VS.B01.P01.C08.S15.SS11.P004
\phantomsection
\label{para:VS.B01.P01.C08.S15.SS11.P004}

The comparison also reveals a difference in
explanatory scope. Potentiality and actuality
describe a relation between a subject and its
possible or realized determinations. Identity
cannot be reduced to either pole of that relation.
A substance may undergo actualization without
that fact alone determining the number, status,
or persistence conditions of its identities.


% =============================================================================
% SS12 — CAUSALITY AND THE ETERNAL ORDER
% =============================================================================

\subsection{Causality, Dependence, and the Eternal Order}
\label{subsec:ch08_s15_ss12}

% Structural ID: VS.B01.P01.C08.S15.SS12.P001
\phantomsection
\label{para:VS.B01.P01.C08.S15.SS12.P001}

Averroes's account of sensible being is inseparable
from his theory of natural causality. Matter and
form do not explain a substance in isolation from
the causal conditions through which it comes to
be and undergoes change. Natural agents actualize
the potentialities of material subjects, while
formal and final causes contribute to the explanation
of determinate natural processes. The causal
structure is therefore part of the explanation
of how substances are generated and maintained
\citep{AristotlePhysics,
AverroesIncoherenceVanDenBergh1954,
BenAhmedPasnauIbnRushdSEP2025}.

% Structural ID: VS.B01.P01.C08.S15.SS12.P002
\phantomsection
\label{para:VS.B01.P01.C08.S15.SS12.P002}

Averroes also defends an Aristotelian account
of the eternity of motion and the world. His
discussion distinguishes essentially ordered
causal series from accidentally ordered series.
The former concern causal dependence in which
the prior causes are required for the ongoing
operation of the series; the latter concern
temporal succession in which earlier stages
may cease while later stages continue. This
distinction is central to his response to
arguments that treat every infinite temporal
series as though it were an infinite simultaneous
causal dependence
\citep{AverroesIncoherenceVanDenBergh1954,
AverroesMetaphysicsBookLambdaGenequand1984}.

% Structural ID: VS.B01.P01.C08.S15.SS12.P003
\phantomsection
\label{para:VS.B01.P01.C08.S15.SS12.P003}

The eternity of the world and the detailed
celestial cosmology belong to Averroes's
historical Aristotelian commitments. They
are not adopted as foundational claims of
Book I. Their relevance here is comparative:
they demonstrate how Averroes connects the
existence and activity of natural substances
to an enduring causal order. A historical
reconstruction must preserve this commitment
without turning it into an unargued premise
of the Philosophy of Viability
\citep{AverroesIncoherenceVanDenBergh1954,
BenAhmedPasnauIbnRushdSEP2025}.

% Structural ID: VS.B01.P01.C08.S15.SS12.P004
\phantomsection
\label{para:VS.B01.P01.C08.S15.SS12.P004}

Causal dependence does not automatically
establish identity. The fact that an existent
depends upon a cause for its generation or
continued activity does not by itself identify
the existent with its cause, with the causal
relation, or with the sequence through which
it is produced. Causality explains relations
of production and dependence; the ontological
status of Identity remains a distinct question.


% =============================================================================
% SS13 — ACCIDENTAL AND SUBSTANTIAL CHANGE
% =============================================================================

\subsection{Accidental Change, Substantial Change, and Continuity}
\label{subsec:ch08_s15_ss13}

% Structural ID: VS.B01.P01.C08.S15.SS13.P001
\phantomsection
\label{para:VS.B01.P01.C08.S15.SS13.P001}

Averroes distinguishes accidental change from
substantial change. In accidental change, a
substance may acquire or lose features while
remaining the subject of the alteration. In
substantial change, by contrast, the departure
of the substantial form marks a deeper
transformation in which the former substance
ceases to exist as that substance, or is
spoken of only homonymously. The distinction
prevents all change from being described
as alteration of one and the same subject
in an identical sense
\citep{AverroesEpitomeMetaphysicsArnzen2010,
BenAhmedPasnauIbnRushdSEP2025}.

% Structural ID: VS.B01.P01.C08.S15.SS13.P002
\phantomsection
\label{para:VS.B01.P01.C08.S15.SS13.P002}

The difference between accidental and substantial
change creates a direct challenge for any
theory that identifies Identity with persistence.
A substance may persist through accidental
alterations, but persistence is not thereby
the definition of Identity. Conversely, the
loss of a substantial form may mark the end
of a substance without establishing a general
rule that every substantial transformation
corresponds to one and only one change of
Identity in every possible ontological context.

% Structural ID: VS.B01.P01.C08.S15.SS13.P003
\phantomsection
\label{para:VS.B01.P01.C08.S15.SS13.P003}

Book I distinguishes Identity from Persistence
and Continuity. These distinctions are not
introduced to deny that persistence and
continuity matter. They are required because
the persistence of a subject, the continuity
of a process, and the Identity of an existent
are not interchangeable explanatory notions.
Averroes's distinction between accidental
and substantial change sharpens the problem
without resolving it in advance.

% Structural ID: VS.B01.P01.C08.S15.SS13.P004
\phantomsection
\label{para:VS.B01.P01.C08.S15.SS13.P004}

The open comparative question is whether
the Aristotelian distinction between substantial
and accidental change can be integrated into
Book I's account of identity-bearing beings
without making substantial form the universal
criterion of Identity. The answer cannot be
assumed from the distinction itself; it
requires a separate analysis of the relation
between determination, transformation,
persistence, and identity.


% =============================================================================
% SS14 — THE STRONGEST AVERROIST CHALLENGE
% =============================================================================

\subsection{The Strongest Averroist Challenge to Book I}
\label{subsec:ch08_s15_ss14}

% Structural ID: VS.B01.P01.C08.S15.SS14.P001
\phantomsection
\label{para:VS.B01.P01.C08.S15.SS14.P001}

The strongest Averroist challenge is that the
concept of Identity may be explanatorily
redundant. If the determinate nature of a
sensible substance is accounted for through
its substantial form, if its material conditions
are explained through matter, if its accidental
features are understood through their dependence
upon substance, and if its generation and
alteration are explained through causal powers,
then an additional ontological category called
Identity may appear unnecessary.

% Structural ID: VS.B01.P01.C08.S15.SS14.P002
\phantomsection
\label{para:VS.B01.P01.C08.S15.SS14.P002}

The challenge can be stated more strongly.
If Identity is not matter, form, essence,
individuation, unity, numerical identity,
structure, persistence, or causal dependence,
what does it contribute to the explanation
of an existent? If it does not distinguish
a substance from another substance, determine
its nature, or account for its capacity to
undergo change, then it may appear to name
a problem already resolved by the Aristotelian
account of substance.

% Structural ID: VS.B01.P01.C08.S15.SS14.P003
\phantomsection
\label{para:VS.B01.P01.C08.S15.SS14.P003}

This challenge is not answered by declaring
that Identity is fundamental. A foundational
commitment can establish the starting point
of an ontology, but it cannot replace the
argument required to show the explanatory
significance of that commitment. Book I
must therefore demonstrate that Identity
is neither an additional component of the
substance nor a redundant name for the
determinate nature explained by substantial
form.

% Structural ID: VS.B01.P01.C08.S15.SS14.P004
\phantomsection
\label{para:VS.B01.P01.C08.S15.SS14.P004}

The Averroist challenge is strongest when
it asks for a positive account rather than
a list of exclusions. It requires Book I
to state what question Identity answers,
what distinctions become possible through
it, and why the account of matter, form,
essence, and causality does not exhaust
that question. The burden of proof therefore
falls on the positive development of Identity,
not on the mere assertion that it cannot be
reduced to another concept.


% =============================================================================
% SS15 — THE BOOK I CHALLENGE AND AN AVERROIST REPLY
% =============================================================================

\subsection{The Book I Challenge and a Possible Averroist Reply}
\label{subsec:ch08_s15_ss15}

% Structural ID: VS.B01.P01.C08.S15.SS15.P001
\phantomsection
\label{para:VS.B01.P01.C08.S15.SS15.P001}

The strongest challenge from Book I begins
by recognizing the explanatory achievements
of Averroes's ontology. Matter and form
explain important principles of sensible
substance; essence explains the determinate
nature of the thing; individuation and
multiplication address questions of numerical
plurality; causal powers explain generation
and change. The challenge is not that these
concepts are inadequate for their own tasks.
It is that their explanatory tasks do not
automatically coincide with the ontological
question of Identity.

% Structural ID: VS.B01.P01.C08.S15.SS15.P002
\phantomsection
\label{para:VS.B01.P01.C08.S15.SS15.P002}

Book I's foundational relation between
Being and Identity consists of three
independent commitments: Analytical
Distinction, Ontological Inseparability,
and Absence of Ontological Priority.
These claims do not introduce two independently
existing entities, nor do they describe
a composition in which Being and Identity
are two material or formal parts. They
state that the two are analytically
distinguishable, ontologically inseparable,
and not ordered by an established priority
of one over the other.

% Structural ID: VS.B01.P01.C08.S15.SS15.P003
\phantomsection
\label{para:VS.B01.P01.C08.S15.SS15.P003}

The distinction allows Book I to ask
what it means for a being to be
identity-bearing without treating Identity
as an additional property or as a cause
that produces Being. It also prevents
the reduction of Identity to the determinate
form, essence, material constitution,
or persistence of an existent. The
ontological relation is not a third
substance added to the Aristotelian
composite; it is a distinct analytical
question concerning the existent itself.

% Structural ID: VS.B01.P01.C08.S15.SS15.P004
\phantomsection
\label{para:VS.B01.P01.C08.S15.SS15.P004}

A possible Averroist reply remains available.
Averroes could argue that if Identity does
not add a distinct constituent, cause,
or explanatory principle to the substance,
then the distinction has no independent
metaphysical work to perform. The Book I
response must therefore show that analytical
distinction does not require an additional
entity and that ontological inseparability
does not make the distinction redundant.
This response is a task for Book I's
positive theory, not a result established
by the comparison alone.


% =============================================================================
% SS16 — BILATERAL EXPLANATORY COSTS AND SUBSTITUTION TEST
% =============================================================================

\subsection{Bilateral Explanatory Costs and the Substitution Test}
\label{subsec:ch08_s15_ss16}

% Structural ID: VS.B01.P01.C08.S15.SS16.P001
\phantomsection
\label{para:VS.B01.P01.C08.S15.SS16.P001}

The comparison can be tested by substituting
Averroist concepts for Identity in the
principal commitments of Book I. If
Identity is replaced by substantial form,
the account gains an explanation of the
determinate nature and essence of the
substance, but it loses the distinction
between Identity and formal determination.
If Identity is replaced by matter, the
account gains a principle of material
potentiality and corporeal multiplication,
but it loses the distinction between
Identity and material constitution.

% Structural ID: VS.B01.P01.C08.S15.SS16.P002
\phantomsection
\label{para:VS.B01.P01.C08.S15.SS16.P002}

If Identity is replaced by individuation
or numerical unity, the account gains
a way to discuss the distinction of
individuals, but it cannot thereby explain
the full relation between Being and Identity.
If Identity is replaced by persistence,
the account gains a temporal criterion
for the continuation of a subject, but
it confuses the continuation of a process
or substance with the ontological status
of Identity. None of these substitutions
preserves all the claims of Book I.

% Structural ID: VS.B01.P01.C08.S15.SS16.P003
\phantomsection
\label{para:VS.B01.P01.C08.S15.SS16.P003}

The substitution test does not prove that
Book I has already established a complete
positive theory of Identity. It shows
that the available Aristotelian concepts
cannot be substituted without changing
the questions and explanatory functions
involved. The result is therefore conditional:
the distinction of Identity remains
philosophically meaningful only if Book I
can develop its positive content and
show that the distinction is not redundant.

% Structural ID: VS.B01.P01.C08.S15.SS16.P004
\phantomsection
\label{para:VS.B01.P01.C08.S15.SS16.P004}

The explanatory cost of Averroes's framework,
as reconstructed here, is not that it lacks
an account of determinate substances.
It is that its central explanatory resources
are organized around substance, matter,
form, essence, and causal actuality rather
than around the independent analytical
question formulated by Book I. The
explanatory cost for Book I is correspondingly
different: it must show why Identity is
needed without converting it into an
additional component of substance.


% =============================================================================
% SS17 — IMPLICATIONS FOR BEING AND IDENTITY
% =============================================================================

\subsection{Implications for Being and Identity}
\label{subsec:ch08_s15_ss17}

% Structural ID: VS.B01.P01.C08.S15.SS17.P001
\phantomsection
\label{para:VS.B01.P01.C08.S15.SS17.P001}

The comparison with Averroes reinforces
the need to preserve the three foundational
claims concerning Being and Identity.
Analytical Distinction permits the two
concepts to be examined separately.
Ontological Inseparability prevents them
from being treated as independently
constituted realities. Absence of
Ontological Priority prevents either
from being established as the ontological
ground of the other. These commitments
must remain distinct because none follows
automatically from the other.

% Structural ID: VS.B01.P01.C08.S15.SS17.P002
\phantomsection
\label{para:VS.B01.P01.C08.S15.SS17.P002}

Averroes's ontology clarifies why the
conceptual boundaries matter. Matter,
form, essence, unity, individuation,
numerical identity, causal dependence,
and persistence are not interchangeable.
Each contributes to a different explanatory
task. Book I must preserve these distinctions
while developing the positive role of
Identity. The result is not a rejection
of substantial determination but a refusal
to treat it as an automatic substitute
for the distinct question of Identity.

% Structural ID: VS.B01.P01.C08.S15.SS17.P003
\phantomsection
\label{para:VS.B01.P01.C08.S15.SS17.P003}

The following diagram places the
Averroist account of sensible substance
beside the foundational relation between
Being and Identity in Book I. The
two panels represent different explanatory
architectures. The relation between
Being and Identity is not depicted as
a composition of matter and form, and
the concepts surrounding it are not
presented as components of Identity.

% Structural ID: VS.B01.P01.C08.S15.SS17.P004
\phantomsection
\label{para:VS.B01.P01.C08.S15.SS17.P004}

% =============================================================================
% Project : THE PHILOSOPHY OF VIABILITY
% Book    : Book I — Ontology of Viability
% Book ID : VS.B01
% Part    : Part I — Identity and Manifestation
% Part ID : VS.B01.P01
% Chapter : Chapter 08 — Comparative Philosophical Examination
% Chapter ID : VS.B01.P01.C08
%
% File :
%   4_ASSETS/diagrams/
%   VS.B01.P01.C08.S15.SS17.P004.DG01.averroes_and_book_i_nonreduction.tex
%
% Asset :
%   DG01 — Averroes's Account of Substance and the Book I
%          Relation between Being and Identity
%
% Structural Identifier :
%   VS.B01.P01.C08.S15.SS17.P004.DG01
%
% Purpose :
%   To distinguish the explanatory architecture of Averroes's
%   sensible substance from the foundational relation between
%   Being and Identity in Book I.
%
% Philosophical Limits :
%   1. The upper structure is a schematic representation of
%      Averroes's account of sensible substance.
%   2. The lower structure is not a composition of two entities.
%   3. The relation between Being and Identity is not represented
%      as a causal or temporal sequence.
%   4. Matter, form, essence, unity, individuation, persistence,
%      and manifestation are not presented as components of Identity.
%   5. The diagram does not claim that Averroes formulated the
%      Book I distinction.
%
% Dependencies :
%   tikz, positioning, arrows.meta.
%   These packages are already loaded by the Book I preamble.
%
% =============================================================================

\begin{figure}[htbp]
\centering

\begin{tikzpicture}[
    font=\small,
    matternode/.style={
        draw=ViabilityBlue,
        fill=ViabilityPaleBlue,
        rounded corners=2pt,
        line width=0.8pt,
        text width=5.0cm,
        align=center,
        inner sep=7pt,
        minimum height=1.1cm
    },
    formnode/.style={
        draw=ViabilityGold,
        fill=ViabilityPaleGold,
        rounded corners=2pt,
        line width=0.8pt,
        text width=5.0cm,
        align=center,
        inner sep=7pt,
        minimum height=1.1cm
    },
    composite/.style={
        draw=ViabilityGreen,
        fill=ViabilityPaleGreen,
        rounded corners=2pt,
        line width=0.9pt,
        text width=10.8cm,
        align=center,
        inner sep=8pt,
        minimum height=1.2cm
    },
    beingnode/.style={
        draw=ViabilityDarkBlue,
        fill=ViabilityPaleBlue,
        rounded corners=2pt,
        line width=0.9pt,
        text width=4.8cm,
        align=center,
        inner sep=8pt,
        minimum height=1.0cm
    },
    identitynode/.style={
        draw=ViabilityGold,
        fill=ViabilityPaleGold,
        rounded corners=2pt,
        line width=0.9pt,
        text width=4.8cm,
        align=center,
        inner sep=8pt,
        minimum height=1.0cm
    },
    note/.style={
        draw=gray,
        fill=white,
        rounded corners=2pt,
        text width=10.8cm,
        align=center,
        inner sep=7pt
    },
    relation/.style={
        -{Stealth[length=2mm]},
        draw=ViabilityDarkBlue,
        line width=0.8pt
    },
    nondirectional/.style={
        draw=ViabilityDarkBlue,
        line width=1.0pt
    },
    boundary/.style={
        densely dashed,
        draw=gray,
        line width=0.75pt
    },
    title/.style={
        font=\bfseries,
        align=center,
        text width=11.0cm
    }
]

% ---------------------------------------------------------------------
% PANEL A — AVERROIST EXPLANATORY ARCHITECTURE
% ---------------------------------------------------------------------

\node[title] (titleA) at (0,0)
    {Averroes: The Constitution of Sensible Substance};

\node[matternode] (matter) at (-3.0,-1.55) {
    \textbf{Prime Matter}\\[3pt]
    Potential subject for substantial forms
};

\node[formnode] (form) at (3.0,-1.55) {
    \textbf{Substantial Form}\\[3pt]
    Actuality and determinate nature
};

\node[composite] (substance) at (0,-3.45) {
    \textbf{Composite Sensible Substance}\\[4pt]
    A determinate existent constituted through matter and form
};

\node[composite] (accidents) at (0,-5.05) {
    \textbf{Accidental Determinations}\\[3pt]
    Dependent features that may be gained or lost while
    the substance remains their subject
};

\draw[relation]
    (matter.south) --
    node[fill=white,inner sep=2pt,font=\scriptsize]
    {material potentiality}
    (substance.north west);

\draw[relation]
    (form.south) --
    node[fill=white,inner sep=2pt,font=\scriptsize]
    {actualization and determination}
    (substance.north east);

\draw[relation]
    (substance) --
    node[fill=white,inner sep=2pt,font=\scriptsize]
    {subject and dependence}
    (accidents);

% ---------------------------------------------------------------------
% PANEL B — BOOK I FOUNDATIONAL RELATION
% ---------------------------------------------------------------------

\node[title] (titleB) at (0,-6.35)
    {Book I: The Relation between Being and Identity};

\node[beingnode] (being) at (-3.0,-7.65) {
    \textbf{Being}
};

\node[identitynode] (identity) at (3.0,-7.65) {
    \textbf{Identity}
};

\draw[nondirectional]
    (being.east) --
    node[
        fill=white,
        text=ViabilityDarkBlue,
        font=\scriptsize,
        align=center,
        inner sep=3pt
    ] {
        Analytical Distinction\\
        Ontological Inseparability\\
        Absence of Ontological Priority
    }
    (identity.west);

\node[note] (nonreduction) at (0,-9.45) {
    \textbf{Non-equivalence}\\[3pt]
    Matter \quad Form \quad Essence \quad Unity \quad
    Individuation \quad Persistence \quad Manifestation
    \\[3pt]
    These concepts perform distinct explanatory functions;
    none is identified with Identity.
};

\draw[boundary]
    (being.south) -- (nonreduction.north west);

\draw[boundary]
    (identity.south) -- (nonreduction.north east);

\node[
    align=center,
    font=\scriptsize,
    text width=10.8cm
] at (0,-10.45) {
    The two panels represent different explanatory questions.
    The lower relation is not a material or formal composition,
    and the dashed boundaries do not indicate causal dependence.
};

\end{tikzpicture}

\caption[
    Averroes's account of substance and the Book I relation
    between Being and Identity
]{
    Schematic comparison between Averroes's account of
    sensible substance through matter and substantial form
    and the foundational relation between Being and Identity
    in Book I. The diagram distinguishes the explanatory
    roles of matter, form, essence, and accidental determination
    from the analytical distinction and ontological inseparability
    of Being and Identity.
}

\label{fig:VS.B01.P01.C08.S15.SS17.P004.DG01}

\end{figure}

% =============================================================================
% END OF DIAGRAM DG01
% =============================================================================


% =============================================================================
% SS18 — OPEN PROBLEMS AND CLAIM STATUS
% =============================================================================

\subsection{Open Problems and Claim Status}
\label{subsec:ch08_s15_ss18}

% Structural ID: VS.B01.P01.C08.S15.SS18.P001
\phantomsection
\label{para:VS.B01.P01.C08.S15.SS18.P001}

The comparison leaves several problems
unresolved. First, Book I must provide
a positive account of Identity that
distinguishes it from essence and
substantial form without treating it
as a separate constituent. Second, the
relation between Identity and the
individuation of concrete substances
requires further analysis. Third, the
relation between substantial change
and the transformation of Identity
cannot be inferred solely from the
Aristotelian distinction between
accidental and substantial change.

% Structural ID: VS.B01.P01.C08.S15.SS18.P002
\phantomsection
\label{para:VS.B01.P01.C08.S15.SS18.P002}

Further questions concern the status
of matter as a potential principle,
the relation between essence and
existence, the unity of a substance,
and the explanatory role of causal
dependence. These issues should not
be converted into settled Book I
principles merely because Averroes
offers a historically influential
account of them. Their comparative
value lies in identifying the precise
questions that the ontology must
continue to address.

% Structural ID: VS.B01.P01.C08.S15.SS18.P003
\phantomsection
\label{para:VS.B01.P01.C08.S15.SS18.P003}

The following table distinguishes
the functional or logical form of
the principal claims from their
epistemic or justificatory status.
Historical reconstruction is not
itself an epistemic status. The
status assigned to a claim indicates
whether it is adopted by Book I,
argued within the comparative
reconstruction, conditionally supported,
provisional, unresolved, or rejected.

% Structural ID: VS.B01.P01.C08.S15.SS18.P004
\phantomsection
\label{para:VS.B01.P01.C08.S15.SS18.P004}

% =============================================================================
% Project : THE PHILOSOPHY OF VIABILITY
% Book    : Book I — Ontology of Viability
% Book ID : VS.B01
% Part    : Part I — Identity and Manifestation
% Part ID : VS.B01.P01
% Chapter : Chapter 08 — Comparative Philosophical Examination
% Chapter ID : VS.B01.P01.C08
%
% File :
%   4_ASSETS/tables/
%   VS.B01.P01.C08.S15.SS18.P004.TL01.averroes_claims_and_comparative_limits.tex
%
% Asset :
%   TL01 — Averroes's Principal Claims and Comparative Limits
%
% Structural Identifier :
%   VS.B01.P01.C08.S15.SS18.P004.TL01
%
% Purpose :
%   To distinguish the functional or logical form of the
%   reconstructed claims from their epistemic or justificatory
%   status, and to record their limits for Book I.
%
% Dependencies :
%   longtable, booktabs, array.
%   These packages are already loaded by the Book I preamble.
%
% =============================================================================

\begingroup
\small
\renewcommand{\arraystretch}{1.18}
\setlength{\tabcolsep}{4pt}
\setlength{\LTcapwidth}{\linewidth}

\def\colw#1{\dimexpr(\linewidth-8\tabcolsep)*#1/100\relax}

\def\rowsep{\specialrule{0.2pt}{3pt}{3pt}}

\def\tablehead{%
    \toprule
    \textbf{Claim or issue} &
    \textbf{Averroist reconstruction} &
    \textbf{Functional or logical form} &
    \textbf{Epistemic or justificatory status} &
    \textbf{Comparative limit for Book I}
    \\
    \midrule
}

\def\grouprow#1{%
    \multicolumn{5}{@{}l}{\footnotesize\textsc{#1}} \\
    \cmidrule{1-5}
    \addlinespace[1pt]
}

\begin{longtable}{@{}
    >{\raggedright\arraybackslash}p{\colw{14}}
    >{\raggedright\arraybackslash}p{\colw{26}}
    >{\raggedright\arraybackslash}p{\colw{14}}
    >{\raggedright\arraybackslash}p{\colw{17}}
    >{\raggedright\arraybackslash}p{\colw{29}}
@{}}

\caption{Averroes's Principal Claims and Their Comparative Limits}
\label{tab:VS.B01.P01.C08.S15.SS18.P004.TL01}
\\

\tablehead
\endfirsthead

\multicolumn{5}{@{}l}{%
    \footnotesize\textit{
        Table~\thetable\ (continued):
        Averroes's Principal Claims and Their Comparative Limits
    }
}
\\
\tablehead
\endhead

\multicolumn{5}{@{}r}{%
    \footnotesize\textit{Continued on next page}
}
\\
\endfoot

\bottomrule
\endlastfoot

\grouprow{Being and substance}

\textbf{Being qua being}
&
Metaphysics investigates being universally and the principles
that belong to existents across the special sciences.
&
Thesis
&
Argued/Derived Claims
&
The subject of metaphysics does not by itself settle
the relation between Being and Identity in Book I.
\\
\rowsep

\textbf{Substance}
&
Substance occupies a primary position among the modes
of being; accidents depend upon substance.
&
Thesis
&
Argued/Derived Claims
&
Substance is not thereby identical with Identity.
\\
\rowsep

\grouprow{Matter, form, and essence}

\textbf{Prime matter}
&
Matter is considered as a potential subject for
substantial forms and natural change.
&
Thesis
&
Argued/Derived Claims
&
Material potentiality is not Identity or Viability.
\\
\rowsep

\textbf{Substantial form}
&
Form actualizes matter and accounts for the determinate
nature of the composite substance.
&
Thesis
&
Argued/Derived Claims
&
Formal determination is not equivalent to Identity.
\\
\rowsep

\textbf{Essence and quiddity}
&
The essence of a sensible substance is explained
through its substantial form.
&
Proposition
&
Conditionally Supported Claims
&
Essence must not be renamed Identity.
\\
\rowsep

\grouprow{Existence and Avicenna}

\textbf{Existence and essence}
&
Averroes criticizes the treatment of existence as
a common attribute or accident of essence.
&
Thesis
&
Conditionally Supported Claims
&
The interpretation of Avicenna's position remains
disputed; the critique is not treated as decisive.
\\
\rowsep

\textbf{Avicennian distinction}
&
Avicenna distinguishes quiddity considered in itself
from the actuality of existence and its causal conditions.
&
Thesis
&
Argued/Derived Claims
&
The distinction is not equivalent to the Being--Identity
distinction.
\\
\rowsep

\grouprow{Unity, individuation, and change}

\textbf{Unity and multiplicity}
&
The one and the many are examined as central
metaphysical problems within the Aristotelian framework.
&
Thesis
&
Argued/Derived Claims
&
Unity is not identical with Identity, Singularity,
or Aggregate.
\\
\rowsep

\textbf{Material multiplication}
&
Extended matter contributes to the multiplication
of material forms.
&
Proposition
&
Conditionally Supported Claims
&
Multiplication is not a universal principle of Identity.
\\
\rowsep

\textbf{Individuation}
&
Forms are individual in their existence; matter
contributes to material plurality.
&
Thesis
&
Conditionally Supported Claims
&
Individuation is not equivalent to Identity.
\\
\rowsep

\textbf{Accidental change}
&
A substance may gain or lose accidents while remaining
the subject of those changes.
&
Proposition
&
Argued/Derived Claims
&
Persistence through accidents is not a definition of Identity.
\\
\rowsep

\textbf{Substantial change}
&
The departure of substantial form marks the end
of the former substance or its continuation only
homonymously.
&
Thesis
&
Conditionally Supported Claims
&
No universal Book I criterion of Identity follows
without further argument.
\\
\rowsep

\grouprow{Causality and cosmology}

\textbf{Natural causality}
&
Natural agents actualize material potentialities
through causal processes.
&
Thesis
&
Argued/Derived Claims
&
Causal dependence does not establish identity with
the cause or the causal process.
\\
\rowsep

\textbf{Eternity of the world}
&
Averroes defends an Aristotelian account of eternal
motion and the world.
&
Thesis
&
Conditionally Supported Claims
&
This is a historical cosmological commitment,
not a foundational commitment of Book I.
\\
\rowsep

\grouprow{Comparative status}

\textbf{The so-called two truths}
&
The inherited description of Averroes as affirming
two contradictory truths is rejected.
&
Rejected Thesis
&
Rejected/Superseded Claims
&
The label is not used as a characterization of
Averroes's position.
\\
\rowsep

\textbf{Being and Identity}
&
Book I maintains Analytical Distinction, Ontological
Inseparability, and Absence of Ontological Priority.
&
Foundational Commitment
&
Foundationally Adopted Commitments
&
These three commitments remain separate and are
not derived from Averroes.
\\
\rowsep

\textbf{Positive role of Identity}
&
The explanatory role of Identity beyond substantial
form, essence, individuation, and persistence
requires further development.
&
Open Problem
&
Unresolved/Open Problems
&
The positive content of Identity must be established
rather than assumed.
\\
\rowsep

\textbf{Comparative synthesis}
&
Averroes's ontology and Book I address overlapping
problems through distinct explanatory architectures.
&
Thesis
&
Provisional Constructions
&
The comparison establishes critical proximity,
not conceptual equivalence.
\\

\end{longtable}

\endgroup

% =============================================================================
% END OF TABLE TL01
% =============================================================================


% =============================================================================
% SS19 — TRANSITION TO IBN ARABI
% =============================================================================

\subsection{Transition to Ibn Arabi}
\label{subsec:ch08_s15_ss19}

% Structural ID: VS.B01.P01.C08.S15.SS19.P001
\phantomsection
\label{para:VS.B01.P01.C08.S15.SS19.P001}

Averroes's analysis has brought the
comparison to the constitution of
sensible substance through matter,
form, essence, potentiality, actuality,
and causal order. The next examination
moves toward a different philosophical
architecture in which the relation
between Being, manifestation, and
ontological disclosure becomes central.
The transition does not establish a
doctrinal continuity between the two
systems. It opens a new comparative
problem.

% Structural ID: VS.B01.P01.C08.S15.SS19.P002
\phantomsection
\label{para:VS.B01.P01.C08.S15.SS19.P002}

The subsequent examination of Ibn Arabi
must therefore determine, on its own
terms, how existence, determination,
manifestation, and multiplicity are
related, without assuming that the
language of substantial form or
Aristotelian causality already provides
the structure of that account.


% =============================================================================
% END OF SECTION 15 OF CHAPTER 08
% =============================================================================